\documentclass[10pt, twocolumn, landscape, a4paper]{article}

\usepackage[margin=1in]{geometry}
\usepackage[utf8]{inputenc}
\usepackage[spanish,es-noshorthands]{babel}
\usepackage{amssymb, amsmath, amsbsy}
\usepackage{mathpazo}
\usepackage{tikz}
% \usetikzlibrary{angles,quotes}
\usepackage[pdftex]{hyperref}
\hypersetup{pdftitle={Elementos gráficos de Geometría},pdfauthor={Luis Carrillo Gutiérrez}}

\renewcommand{\familydefault}{\sfdefault}
\pagestyle{empty} 

\begin{document}
\subsubsection*{Punto}
\begin{tikzpicture}
	\draw (-1, 1) -- (2.75, -1.8);
	\draw (-1, -1) -- (3.5, 2.4);
	\draw (.325, 0) circle (.05);
	\draw (.325, -.25) node {P};
\end{tikzpicture}

\subsubsection*{Línea recta}

\begin{tikzpicture}
	\draw (-3, 0) -- (3.75, 0);
	\draw (.25, -.25) node {r};
\end{tikzpicture}

\subsubsection*{Semirrecta}

\begin{tikzpicture}
	\draw (0, 0) -- (3.75, 0);
	\draw (0, 0) circle (.05);
	\draw (.25, -.25) node {O};
	\draw (1.75, .25) node {r};
\end{tikzpicture}

\subsubsection*{Segmento}

\begin{tikzpicture}
	\draw (0, 0) -- (4.25, 0);
	\draw (0, 0) circle (.05);
	\draw (4.25, 0) circle (.05);
	\draw (.25, -.25) node {A};
	\draw (4, -.25) node {B};
\end{tikzpicture}

\end{document}
